\documentclass[12pt,a4paper]{article}
\usepackage[UTF8]{ctex}
\usepackage{amsmath}
\usepackage{amsfonts}
\usepackage{amssymb}
\usepackage{graphicx}
\usepackage{booktabs}
\usepackage{array}
\usepackage{geometry}
\geometry{left=2.5cm,right=2.5cm,top=2.5cm,bottom=2.5cm}
\usepackage{float}
\usepackage{longtable}

\title{实验 4.9 空气中声速的测定}
\author{}
\date{}

\begin{document}

\maketitle

\subsection*{引言}

声波是一种在弹性介质中传播的机械纵波。频率在 20 Hz ~ 20 kHz 的声波可以被人耳听到，称为可闻声波；频率低于 20 Hz 的声波称为次声波；频率高于 20 kHz 的声波称为超声波。次声波和超声波都是人耳不能听到和辨别的。在同一介质中，声速基本与频率无关，所以超声波的传播速度就是声速。由于超声波具有波长短、易于定向发射、不会造成听觉污染等特点，因此本实验采用超声波测量声速。实际应用中，超声波在测距、定位、测液体的流速、测比重、测溶液的浓度、测材料的杨氏模量、考察气体温度变化等方面都有重要意义。

\subsection*{实验目的}

\begin{enumerate}
\item 熟练掌握用共振干涉法（驻波法）、相位法和时差法测量声波在空气中的传播速度。
\item 进一步熟悉示波器的使用。
\item 学会运用逐差法处理测量数据。
\end{enumerate}

\subsection*{实验仪器}

FD-SV-A 型声速测量实验仪、（双踪）示波器。

\subsection*{实验原理}

\begin{enumerate}
\item \textbf{超声波的产生与接收}

超声波的产生与接收可以由两只结构完全相同的超声压电陶瓷换能器分别完成，压电陶瓷换能器可以实现声压和电压之间的转换。超声波由压电陶瓷换能器将电信号通过逆压电效应（电声转换）产生并发射；超声波也可由压电陶瓷换能器接收并通过压电效应（声电转换）再转换成电信号。压电陶瓷换能器在声电转换过程中保持信号频率不变。如图 4-9-1 所示，信号源将几万赫兹的交流电信号送入压电陶瓷换能器 $S_1$，$S_1$ 由逆压电效应发出一平面超声波，传输给压电陶瓷换能器 $S_2$，$S_2$ 由压电效应将接收到的声压转换成电信号，送入示波器后，通过示波器便可观察到电信号的大小强弱变化。

压电陶瓷换能器具有固有的谐振频率，当输入电信号的频率等于谐振频率时，压电陶瓷换能器处于谐振状态，此时它的振幅最大，发射器发出超声波的功率最大，仪器处于最佳工作状态。

\item \textbf{测量声速的实验方法}

声速的测量方法可以分为两大类：一是直接法（时差法），即测读超声波的传输时间 $t$ 及距离 $L$，由 $v=\frac{L}{t}$ 求出声速 $v$；二是间接法（波长 - 频率法），利用关系 $v=\lambda f$，测出其频率 $f$ 和波长 $\lambda$，即可求出声速 $v$。本实验采用的共振干涉法（驻波法）和相位法就属于间接法。

\begin{itemize}
\item \textbf{（1）共振干涉法测声速}

如图 4-9-1 所示，由超声压电陶瓷换能器（发射器）$S_1$ 发出的超声波，经介质（空气）传播到超声压电陶瓷换能器（接收器）$S_2$，$S_2$ 在接收超声波信号的同时还反射一部分超声波信号。如果接收器 $S_2$ 与发射器 $S_1$ 严格平行，则发射波与反射波沿相反方向传播，振动方向一致，且振幅大致相等，两者满足相干条件，相干叠加即形成驻波。

假设发射波和反射波的波函数分别为：
\begin{equation}
y_1=A\cos(\omega t-\frac{2\pi}{\lambda}x), \quad (4-9-1)
\end{equation}
\begin{equation}
y_2=A\cos(\omega t+\frac{2\pi}{\lambda}x), \quad (4-9-2)
\end{equation}
则两列波叠加后合成波的波函数为：
\begin{equation}
y=y_1+y_2=2A\cos\frac{2\pi}{\lambda}x\cos\omega t. \quad (4-9-3)
\end{equation}
当 $\cos\frac{2\pi}{\lambda}x=1$，即 $\frac{2\pi}{\lambda}x=n\pi$（$n=0,1,2,\cdots$）时，振幅最大，即当两个超声压电陶瓷换能器之间的距离等于半波长的整数倍时，将产生驻波现象。

根据波的干涉理论，接收器两相邻声压极大值之间的距离为 $\frac{\lambda}{2}$。因此，若保持频率 $f$ 不变，通过测量两相邻信号幅值最大时 $S_1$ 与 $S_2$ 之间的距离差，则可以得到声波的半波长 $\frac{\lambda}{2}$，由 $v=\lambda f$ 即可求出声速 $v$。

\item \textbf{（2）相位法测声速}

如果用示波器测李萨如图形的方法观察超声波的发射波形与接收波形，则屏幕上一般会显示一个倾斜的椭圆，这是由于发射波在空间传播一段距离后，会有相位变化，椭圆倾斜情况反映发射与接收信号之间的相位差。缓慢改变接收器的位置，接收信号的相位会随之缓慢变化，导致椭圆倾斜角度随之变化，如图 4-9-3 所示。当椭圆变成通过第一、三象限的一条直线段时，发射信号与接收信号同相，相位差为 0；继续移动接收器，当椭圆变成通过第二、四象限的一条直线段时，发射信号与接收信号反相，相位差为 $\pi$。相位差从 0 到 $\pi$，接收器移动的距离应为半个波长，因此可以精确测定超声波在空气中传播时的波长，由 $v=\lambda f$ 即可求出声速 $v$。

\item \textbf{（3）时差法测声速}

将施加在发射器 $S_1$ 上的连续正弦电信号换成一脉冲电信号，超声发射器 $S_1$ 产生一在介质中传播的超声脉冲，经过时间 $t$ 后，到达距离 $L$ 处的接收器 $S_2$，测读超声波的传输时间 $t$ 及距离 $L$，由 $v=\frac{L}{t}$ 可求出声速 $v$。实际测量时，通常采用测量距离变化 $\Delta L$ 和时间变化 $\Delta t$ 的方法求出声速 $v=\frac{\Delta L}{\Delta t}$。
\end{itemize}

\item \textbf{理想气体中的声速值}

声波在理想气体中的传播过程是绝热过程，声速为：
\begin{equation}
v=\sqrt{\frac{\gamma R\Theta}{\mu}}, \quad (4-9-4)
\end{equation}
其中，$\gamma$ 为气体的绝热系数（空气的定压热容与定容热容之比），$R$ 为普适气体常量，$\mu$ 为相对分子质量，$\Theta$ 为气体的热力学温度。由此可见，气体中的声速 $v$ 不仅与温度 $\Theta$ 有关，还与绝热系数 $\gamma$ 和相对分子质量 $\mu$ 有关（$\gamma$ 和 $\mu$ 与气体的成分有关）。

若以摄氏温度 $\theta$ 计算，则 $\Theta=273.15+\theta$。将其代入式 (4-9-4) 可得：
\begin{equation}
v=\sqrt{\frac{\gamma R}{\mu}(273.15+\theta)}=\sqrt{\frac{\gamma R}{\mu}273.15}\cdot\sqrt{1+\frac{\theta}{273.15}}=v_0\sqrt{1+\frac{\theta}{273.15}}. \quad (4-9-5)
\end{equation}
在标准状态下，对于空气，0°C 时的声速为 $v_0=331.45\,\text{m/s}$。若同时考虑到空气中蒸汽的影响，则校准后声速的表达式为：
\begin{equation}
v=331.45\sqrt{1+\frac{\theta}{\Theta_0}}\left(1+\frac{0.319p_w}{p}\right)\,\text{m/s}, \quad (4-9-6)
\end{equation}
其中，$p_w$ 为蒸汽的分压强，$p$ 为大气压强。
\end{enumerate}

\subsection*{仪器介绍}

FD-SV-A 型声速测量实验仪采用小型高效超声波传感器，该传感器有金属外壳屏蔽，抗干扰能力强。长度的测量采用 50 分度的游标卡尺，精度为 0.02 mm。游标卡尺的游标上装有微调螺旋，可以精细调节 $S_2$ 的位置。超声波传感器有两个信号输出端口，输出正弦信号的频率范围为 38 ~ 42 kHz。超声压电陶瓷换能器的振荡频率为 $(40.1\pm0.4)\,\text{kHz}$。

双踪示波器由垂直放大系统、水平扫描与放大系统、示波管显示系统和电源组成。其中，垂直放大系统有调节信号的 "伏 / 格（VOLTS/DIV）" 选择功能，用于控制荧光屏上所显示的 Y 信号的大小；水平扫描与放大系统具有 "时间 / 格（TIME/DIV）" 选择和 "伏 / 格（VOLTS/DIV）" 选择功能，用于控制扫描信号频率或荧光屏上所显示的 X 信号的大小；示波管显示系统电路提供示波管各极的电压，这里，"聚焦（FOCUS）" 和 "辉度（INTEN）" 功能用于控制荧光屏上光点的大小和亮度，而 "◆◆POSITION" 选择和 "▲▼POSITION" 选择功能分别用于控制荧光屏上光点的水平位置和垂直位置。

\subsection*{实验内容}

\begin{enumerate}
\item \textbf{共振干涉法测声速}

\begin{enumerate}
\item 将声速测量实验仪的一个发射信号输出端口与发射器 $S_1$ 相连，接收器 $S_2$ 与示波器的 "CH2" 输入端口相连。
\item 接通示波器，将工作模式（MODE）拨到 "CH2" 处，触发源选择 "CH2"，触发方式选择 "AUTO"，垂直轴输入信号的输入方式选择 "AC"，调节 "时间 / 格（TIME/DIV）" 旋钮，使屏上出现信号波形，调节 "LEVEL" 旋钮，使屏上出现稳定的信号波形。
\item 打开声速测量实验仪开关。将接收器 $S_2$ 相对于发射器 $S_1$ 移开一定距离。调节信号源的频率，当频率接近 40 kHz 左右时缓慢调节，观察示波器的电压幅度，使之达到最大，此时信号源的输出频率即等于换能器系统的谐振频率。记录频率值 $f$，此后整个实验过程中不必再改变信号频率。
\item 调节示波器的 "时间 / 格（TIME/DIV）" 旋钮至屏上出现多个周期波形，调节 "伏 / 格（VOLTS/DIV）" 旋钮使屏上的波形幅值大小适中。
\item 使 $S_1$ 距 $S_2$ 大约 30 mm，缓慢移动 $S_2$，当信号波形幅值达到极大值时即可确认 $S_1$、$S_2$ 间形成驻波。旋转游标卡尺游标后面的微调螺母，精确定位此时 $S_2$ 的位置坐标值 $x_1$。然后继续往大调节 $S_2$ 的位置，使信号波形幅值达到又一极大值时，采用前述方法精确定位 $S_2$ 的位置坐标值 $x_2$。连续精确定位出 10 个极大值时的 $S_2$ 位置坐标值，由逐差法可得：
\begin{equation}
\lambda=\frac{2}{5}(|x_{10}-x_5|+\cdots+|x_6-x_1|). \quad (4-9-7)
\end{equation}
由式 (4-9-7) 即可求出空气中的声波波长。
\item 由 $v=\lambda f$ 求出空气中的声速。
\end{enumerate}

\item \textbf{相位法测声速}

\begin{enumerate}
\item 将声速测量实验仪的另一个发射信号输出端口与示波器的 "CH1" 输入端口相连，接收器 $S_2$ 与示波器的 "CH2" 输入端口相连。
\item 逆时针旋转示波器 "时间 / 格（TIME/DIV）" 旋钮到底部 X=Y 位置，示波器上显示一椭圆，调节 CH1 和 CH2 的 "伏 / 格（VOLTS/DIV）" 旋钮，使椭圆大小合适。
\item 改变 $S_2$ 的位置，观察示波器显示的椭圆变化。椭圆变为 "第一、三象限" 和 "第二、四象限" 的一条直线段时，旋转游标卡尺游标后面的微调螺母，精确定位图形变为直线段时 $S_2$ 的位置坐标值，位置坐标值之差的绝对值即为半个波长。连续精确定位出 10 个椭圆变为直线段时 $S_2$ 的位置坐标值，利用逐差法计算波长及声速。
\end{enumerate}
\end{enumerate}

\end{document}